\section{常微分方程 III}

\subsection{高阶线性微分方程解的结构}

对于一个线性微分方程
\begin{equation} \label{eq3.1}
    y^{(n)} + a_{n - 1}(x) y^{(n - 1)} + \dots + a_1(x) y' + a_0(x) y = f(x)
\end{equation}
我们先考虑其齐次形式
\begin{equation} \label{eq3.2}
    y^{(n)} + a_{n - 1}(x) y^{(n - 1)} + \dots + a_1(x) y' + a_0(x) y = 0
\end{equation}

那么根据线性性，不难得知对于方程 $\ref{eq3.2}$ 的任意两个解，它们的线性组合也是方程 $\ref{eq3.2}$ 的解。而对于 $n$ 阶线性微分方程，其具有 $n$ 个线性无关的特解
$$
y_1(x), y_2(x), \dots, y_n(x)
$$
那么方程 $\ref{eq3.2}$ 的通解为
$$
y(x) = C_1 y_1(x) + C_2 y_2(x) + \dots + C_n y_n(x)
$$

而假设方程 $\ref{eq3.1}$ 有一个特解 $y^*(x)$，那么方程 $\ref{eq3.1}$ 的通解为
$$
y(x) = y^*(x) + C_1 y_1(x) + C_2 y_2(x) + \dots + C_n y_n(x)
$$

\subsection{降阶法}

对于求齐次线性微分方程的特解可以使用降阶法。

已知二阶齐次线性微分方程
\begin{equation} \label{eq3.3}
    y'' + P(x) y' + Q(x) y = 0
\end{equation}
的一个非零特解 $y_1(x)$。如果要求它的另一个线性无关的解 $y_2(x)$，可以令：
$$
y_2(x) = u(x) y_1(x)
$$
将其代入原方程中，这可使得该二阶方程关于 $u'$ 变为了一个一阶方程，从而实现“降阶”的目的。具体过程如下：

对其求导得到：
$$
\begin{aligned}
    y_2' & = u' y_1 + u y_1' \\
    y_2'' & = u'' y_1 + 2 u' y_1' + u y_1''
\end{aligned}
$$
代入原方程并在各项中提取出 $u, u', u''$，整理得到：
$$
u(y_1'' + P(x) y_1' + Q(x) y_1) + u'(2y_1' + P(x) y_1) + u'' y_1 = 0
$$
因为 $y_1(x)$ 是原方程的特解，所以有 $y_1'' + P(x) y_1' + Q(x) y_1 = 0$。于是方程化简为：
$$
u'' y_1 + u' (2y_1' + P(x) y_1) = 0
$$
令 $v = u'$，则 $v' = u''$，上式变为关于 $v$ 的一阶可分离变量微分方程：
$$
v' y_1 + v (2y_1' + P(x) y_1) = 0
$$
移项并分离变量（假设 $y_1 \neq 0$）：
$$
\frac{v'}{v} = - \left(2\frac{y_1'}{y_1} + P(x)\right)
$$
两边对 $x$ 积分得：
$$
\ln \abs{v} = -2 \ln \abs{y_1} - \int P(x) \dd{x} + C
$$
取特解使其最简，即取积分常数 $C = 0$ 且取正号脱去绝对值，得到：
$$
v = u' = \frac{1}{y_1^2} \E^{-\int P(x) \dd{x}}
$$
再进行一次积分，即可求出 $u(x)$：
$$
u(x) = \int \frac{1}{y_1^2} \E^{-\int P(x) \dd{x}} \dd{x}
$$
进而得到了原方程的另一个线性无关解：
$$
y_2(x) = y_1(x) \int \frac{1}{y_1^2} \E^{-\int P(x) \dd{x}} \dd{x}
$$
这样，我们就通过将原方程降一阶，成功求解出了方程的基础解系。

\subsection{二阶常系数线性微分方程的解法}

先考虑下面的二阶常系数齐次线性微分方程：
$$
y'' + py' + qy = 0
$$

我们设 $y = \E^{r x}$，代入得到：
$$
(r^2 + p r + q) \E^{r x} = 0
$$
该方程可以解出两个解 $r_1, r_2$，若：
\begin{itemize}
    \item 两个不同实根 $r_1 = \alpha, r_2 = \beta$：那么原方程通解为
    $$
    y(x) = C_1 \E^{\alpha x} + C_2 \E^{\beta x}
    $$
    \item 两个相同实根 $r_1 = r_2 = \alpha$：那么原方程通解为
    $$
    y(x) = (C_1 + C_2 x) \E^{\alpha x}
    $$
    \item 两个共轭复根 $r_1 = \alpha + \beta i, r_2 = \alpha - \beta i$：那么原方程通解为
    $$
    y(x) = \E^{\alpha x} (C_1 \cos \beta x + C_2 \sin \beta x)
    $$
\end{itemize}

现在需要求方程
$$
y'' + p y' + q y = f(x)
$$
的一个特解。我们假设其对应的齐次方程的通解为 $y = C_1 y_1 + C_2 y_2$，那么使用常数变易法，我们假设原方程的一个特解为
$$
y = C_1(x) y_1 + C_2(x) y_2
$$
代入得到：
$$
y' = C_1'(x) y_1 + C_1(x) y_1' + C_2'(x) y_2 + C_2(x) y_2'
$$
假设有
$$
C_1'(x) y_1 + C_2'(x) y_2 = 0
$$
那么求二阶导：
$$
y'' = C_1(x) y_1'' + C_1'(x) y_1' + C_2(x) y_2'' + C_2'(x) y_2'
$$
代入原微分方程：
$$
\begin{gathered}
    \ab(C_1(x) y_1'' + C_1'(x) y_1' + C_2(x) y_2'' + C_2'(x) y_2') + p \ab(C_1(x) y_1' + C_2(x) y_2') + q \ab(C_1(x) y_1 + C_2(x) y_2) = f(x) \\
    C_1(x) \ab(y_1'' + p y_1' + q y_1) + C_2(x) \ab(y_2'' + p y_2' + q y_2) + C_1'(x) y_1' + C_2'(x) y_2' = f(x) \\
    C_1'(x) y_1' + C_2'(x) y_2' = f(x)
\end{gathered}
$$
结合两个方程组可解得 $C_1(x), C_2(x)$ 的值。

\subsection{作业}

\begin{problem}
    习题 10.3.1
    \begin{solution}
        \item[\textbf{1)}] 先求 $y'' + y' + y = 0$ 的通解：
        $$
        r^2 + r + 1 = 0 \implies r_1 = \frac{-1 + \sqrt{3} \I}{2}, r_2 = \frac{-1 - \sqrt{3} \I}{2}
        $$
        因此该齐次方程通解为
        $$
        y = \E^{-\frac{1}{2} x} \ab(C_1 \cos \frac{\sqrt{3}}{2} x + C_2 \sin \frac{\sqrt{3}}{2} x)
        $$

        而观察得到原方程有一个特解 $y = \frac{1}{3} \E^{x}$。因此原方程通解为：
        $$
        y = \E^{-\frac{1}{2} x} \ab(\cos \frac{\sqrt{3}}{2} x + \sin \frac{\sqrt{3}}{2} x) + \frac{1}{3} \E^{x}
        $$

        \item[\textbf{2)}] 用特征方程法求 $y'' - 3y' + 2y = 0$ 的通解：
        $$
        r^2 - 5r + 6 = 0 \implies r_1 = 2, r_2 = 3
        $$
        因此该齐次方程通解为
        $$
        y = C_1 \E^{2x} + C_2 \E^{3x}
        $$
        假设原方程有一个特解 $y = (a x + b) \E^{4 x}$，那么：
        $$
        y'' - 5 y' + 6y = (x + 1) \E^{4x} \implies 
        \begin{dcases}
            a = \frac{1}{2} \\
            b = -\frac{1}{4}
        \end{dcases}
        $$
        因此原方程通解为
        $$
        y = C_1 \E^{2x} + C_2 \E^{3x} + \ab(\frac{1}{2} x - \frac{1}{4}) \E^{4x}
        $$

        \item[\textbf{3)}] 用特征方程法求 $y'' - 3y' + 2y = 0$ 的通解：
        $$
        r^2 - 3r + 2 = 0 \implies r_1 = 1, r_2 = 2
        $$
        因此该齐次方程通解为
        $$
        y = C_1 \E^{x} + C_2 \E^{2x}
        $$
        假设原方程有一个特解 $y = x (a x + b) \E^{x}$，那么：
        $$
        y'' - 3 y' + 2y = x \E^{x} \implies \begin{dcases}
            a = -\frac{1}{2} \\
            b = -1
        \end{dcases}
        $$
        因此原方程通解为
        $$
        y = C_1 \E^{x} + C_2 \E^{2x} + \ab(-\frac{1}{2} x^2 - x) \E^{x}
        $$

        \item[\textbf{4)}] 用特征方程法求 $y'' + 6y' + 9y = 0$ 的通解：
        $$
        r^2 + 6r + 9 = 0 \implies r_1 = r_2 = -3
        $$
        因此该齐次方程通解为
        $$
        y = (C_1 + C_2 x) \E^{-3x}
        $$
        假设原方程有一个特解 $y = x^2 (a x + b) \E^{-3x}$，那么：
        $$
        y'' + 6y' + 9y = 5x \E^{-3x} \implies \begin{dcases}
            a = \frac{5}{6} \\
            b = 0
        \end{dcases}
        $$
        因此原方程通解为
        $$
        y = (C_1 + C_2 x) \E^{-3x} + \frac{5}{6} x^3 \E^{-3x}
        $$

        \item[\textbf{5)}] 用特征方程法求 $y'' - y' = 0$ 的通解：
        $$
        r^2 - r = 0 \implies r_1 = 0, r_2 = 1
        $$
        因此该齐次方程通解为
        $$
        y = C_1 + C_2 \E^{x}
        $$
        假设原方程有一个特解 $y = (a \sin x + b \cos x) \E^{-x}$，那么：
        $$
        y'' - y' = \E^{-x} \sin x \implies \begin{dcases}
            a = \frac{1}{10} \\
            b = \frac{3}{10}
        \end{dcases}
        $$
        因此原方程通解为
        $$
        y = C_1 + C_2 \E^{x} + \ab(\frac{1}{10} \sin x + \frac{3}{10} \cos x) \E^{-x}
        $$

        \item[\textbf{6)}] 用特征方程法求 $y'' + y = 0$ 的通解：
        $$
        r^2 + 1 = 0 \implies r_1 = \I, r_2 = -\I
        $$
        因此该齐次方程通解为
        $$
        y = C_1 \cos x + C_2 \sin x
        $$
        化简原方程：
        $$
        y'' + y = \cos^2 x = \frac{1 + \cos 2x}{2}
        $$
        可观察得到一个特解
        $$
        y = -\frac{1}{6} \cos 2x + \frac{1}{2}
        $$
        因此原方程通解为
        $$
        y = C_1 \cos x + C_2 \sin x - \frac{1}{6} \cos 2x + \frac{1}{2}
        $$

        \item[\textbf{7)}] 用特征方程法求 $y'' + y = 0$ 的通解：
        $$
        r^2 + 1 = 0 \implies r_1 = \I, r_2 = -\I
        $$
        因此该齐次方程通解为
        $$
        y = C_1 \cos x + C_2 \sin x
        $$
        观察得到原方程有一个特解
        $$
        y = \frac{1}{5} \E^{2x} - \frac{1}{2} x \cos x
        $$
        因此原方程通解为
        $$
        y = C_1 \cos x + C_2 \sin x + \frac{1}{5} \E^{2x} - \frac{1}{2} x \cos x
        $$
    \end{solution}
\end{problem}